\documentclass{beamer}

\newcommand{\darkLogo}{$\vcenter{\includegraphics{tallogo-dark.pdf}}$}

\usefonttheme{serif}
\usepackage[no-math]{fontspec}
\usepackage{amsmath,amsfonts,amssymb,stmaryrd}

\def\yo{\text{y}}
\def\pto{\rightsquigarrow}
\def\Prof{\textbf{Prof}}
\def\Lan{\text{Lan}}

\definecolor{bleakYellow}{HTML}{F0DB4F}
\setmainfont{Source Sans Pro}

\usetheme[background=dark]{metropolis}
\setbeamercolor{alerted text}{fg=bleakYellow}

\setbeamertemplate{footline}{}

\author{Fosco Loregian \darkLogo}
\title{XYZ}
\date{\today}

\setbeamertemplate{navigation symbols}{}

\usepackage[all,2cell]{xy}\UseAllTwocells

\begin{document}
\begin{frame}
  \begin{itemize}
    \item Ph.D. at SISSA - Trieste \includegraphics{ita.pdf}

    \indent {\footnotesize Stable homotopy theory, ∞-categories, derived AG}
    \item University of Western Ontario \includegraphics[scale=.04]{can.pdf}
    
    \indent {\footnotesize ∞-categories, derivators}
    \item Masaryk University \includegraphics{czr.pdf}
    
    \indent {\footnotesize Accessible categories, derivators, 2-categories}
    \item Max Planck Inst. für Math. \includegraphics{ger.pdf}
    
    \indent {\footnotesize 2-categories, derivators, applied category theory}
    \item Centro de Matemàtica da Universidade de Coimbra \includegraphics[scale=.04]{por.pdf}
    
    \indent {\footnotesize 2-categories; finishing my first book}
    \item Tallinna Tehnikaülikooli \includegraphics{est.pdf}
    
    \indent {\footnotesize 2-categories; functorial semantics; categorical probability theory and its applications}
  \end{itemize}
\end{frame}
%
%
%
%
%
\begin{frame}
  \Huge\centering \bfseries STABLE HOMOTOPY THEORY
\end{frame}
%
\begin{frame}
  \alert{∞-categories}: a thickening of the notion of category, suitable for homotopy-invariant mathematics (algebraic geometry, algebraic topology type theory).
  
  Turns out some parts of Mathematics are easier if stated in these terms:
  \begin{itemize}
    \item homological algebra : the scary part of algebraic topology 
    
    {\footnotesize higher algebra: the linear algebra of ∞-categories}
    \item 1-topos theory : a synthetic type theory,
    
    {\footnotesize ∞-topos theory: a synthetic homotopy theory of homotopy types}
  \end{itemize}
\end{frame}
%
%
%
%
%
\begin{frame}
  A stable ∞-category is an ∞-category

with all finite limits and colimits
such that a square is cartesian iff cocartesian

The homotopy category of a stable ∞-cat is always triangulated

The correspondence sending an abelian A into its derived category has a nice and clear universal property

Stable, rational, p-adic, ...  homotopy theory are pieces of commutative algebra of  ∞-categories
\end{frame}
%
%
%
%
%
\begin{frame}
  A t-structure on D triangulated is a pair of triangulated subcategories of D such that every object X lies in a sequence
\[...\]
[FL14] : On stable ∞-cats a t-structure is a factorization system (E,M)
\begin{itemize}
\item such that E and M are 3-for-2 classes
\item thus the category of E-cofibrant objects is coreflective
\item and the category of M-fibrant objects is reflective
\item cof/fib replacement are neg/pos truncation
\end{itemize}
(Algebraic) geometry reduces to a piece of categorical algebra.
\end{frame}
%
%
%
\begin{frame}{Plan: redo $t$-structures}
[FL15] On the set of t-structures there is a natural choice of Z-action (Z = the integers); study Z-equivariant monotone maps from a poset J to TS(C).

Bridgeland stability manifolds, Postnikov towers on ∞-toposes

[FL15b] Every stratified manifold (X,s) generates a pair of t-structure that can be glued together

recollements, stratified schemes, representation of algebras

Project: algebraic geometry is easier, with factorization systems.

Rational and p-adic homotopy ~ rational and p-adic geometry ~ number-theoretic factorization systems

\end{frame}
%
%
%
%
%
\begin{frame}{Every PhD finishes with two conjectures\dots}
 \[\{ J \colon \mathbb R \to TS(\mathcal C) \mid J \text{ is Sorgenfrey continuous} \}\] 
 is an interesting set [FL-PhD]

 \[\{ J \colon \text{Spec}(\mathbb Z) \to TS(\mathcal C) \mid J \text{ is Zariski continuous} \}\]

 Blakers-Massey in char p is a thm about factorization systems
\end{frame}
%
%
%
%
\begin{frame}
  Fact: Bord(n) is the free (∞,n)-symmoncat on the point 

  monoidal functors Z : Bord(n) → Vect are completely classified and used in QFT

  Morse theory is the theory of suitable factorization systems on Bord(n)

  {critical points of a Morse function}

  correspond to

 {critical values of a certain "slicing" J : R → FS(Bord(n))}

 disegno del toro
\end{frame}
%
%
%
%
%
%
%
\begin{frame}
  \Huge\centering \bfseries DERIVATORS
\end{frame}
%
\begin{frame}
  A derivator is a strict 2-functor
  \[\mathbb{D} : {\bf Cat}^\text{op} \to {\bf CAT} \]
  satisfying sheafy conditions. They form the 2-category Der.

  They subsume most of ∞-category theory; in particular, their stable homotopy.
  \begin{itemize}
    \item [LV17] : A t-structure on a stable derivator is still a certain kind of factorization system*

    *And FS are still algebras for the "squaring" 2-monad ( \_ )² : A ↦ A² (see [KT93])
    
    \item  [Lor18] : reflective subderivators correspond to reflective factorization systems, and to algebras for idempotent monads

    (the "formal theory of monads" still holds in Der, a monad T : D → D is just defined objectwise)
  \end{itemize}
\end{frame}
%
\begin{frame}
  There is a Yoneda structure on the 2-category of derivators

  If this is true
\begin{itemize}
\item comprehensive account of various notions of accessible and locally presentable derivator using [DLL18]; categorical logic for derivators; derivator toposes?
\item a convincing form of adjoint functor theorem for derivators: useful, innit?
\item profunctors between derivators; fibered derivators; operads in derivator theory; applications in representation theory, stable homotopy, ...
\end{itemize}
\end{frame}
%
%
%
%
%
\begin{frame}
  \Huge\centering \bfseries FORMAL CATEGORY THEORY
\end{frame}
%
\begin{frame}
  
\end{frame}
%
%
%
%
%
\begin{frame}
  \Huge\centering \bfseries 2-CATEGORIES IN FUNCTIONAL PROGRAMMING
\end{frame}
%
\begin{frame}
  
\end{frame}
%
%
%
%
%
\begin{frame}
  \Huge\centering \bfseries 2-CATEGORIES IN PROBABILITY THEORY
\end{frame}
%
\begin{frame}
  
\end{frame}
%
%
%
%
%
\end{document}